\documentclass[submit]{../harvardml}
% \documentclass[submit]
% \course{CS1810-S25}
\assignment{Homework \#4}
\duedate{April 3, 2026 at 11:59pm}

\usepackage[OT1]{fontenc}
\usepackage[colorlinks,citecolor=blue,urlcolor=blue]{hyperref}
\usepackage[pdftex]{graphicx}
\usepackage{caption}
\usepackage{fullpage}
\usepackage{soul}
\usepackage{amsmath}
\usepackage{amssymb}
\usepackage{framed}
\usepackage{color}
\usepackage{todonotes}
\usepackage{listings}
\usepackage{../common}
\usepackage{float,booktabs}
\newcommand{\E}{\mathbb{E}}
\newcommand{\KL}{D_{\mathrm{KL}}}
\newcommand{\R}{\mathbb{R}}
\newcommand{\bx}{\mathbf{x}}
\newcommand{\bz}{\mathbf{z}}

%%%%%%%%%%%%%%%%%%%%%%%%%%%%%%%%%%%%%%%%%%%
%% Solution environment
\usepackage{xcolor}
\newenvironment{solution}{
    \vspace{2mm}
    \color{blue}\noindent\textbf{Solution}:
}{}
%%%%%%%%%%%%%%%%%%%%%%%%%%%%%%%%%%%%%%%%%%%

\usepackage[mmddyyyy,hhmmss]{datetime}

\definecolor{verbgray}{gray}{0.9}

\lstnewenvironment{csv}{
  \lstset{backgroundcolor=\color{verbgray},
  frame=single,
  framerule=0pt,
  basicstyle=\ttfamily,
  columns=fullflexible}}{}
 
\begin{document}

\begin{center}
{\Large Representation Learning, Transformers, and Non-parametric methods}\\
\end{center}

\subsection*{Introduction}

This homework contains three problems:
\begin{enumerate}
    \item You will do some transformer computations by hand, and implement self-attention from scratch and produce attention score heatmaps.
    \item You will do some derivations related to autoencoders and implement them.
    \item You will do some math to prove bounds on decision tree ensembles, in order to 
\end{enumerate}

We note that this homework does not have any problem where you directly run through the computations you would do to derive a decision tree. Doing so can be good review, however, so for an example with those computations, please see the last exercise in Section 6 notes.

\subsection*{Resources and Submission Instructions}

Please type your solutions after the corresponding problems using this \LaTeX\ template, and start each main problem on a new page.

Please submit the writeup PDF to the Gradescope assignment `HW4'. Remember to assign pages for each question.  \textbf{You must include any plots in your writeup PDF. }. Please submit your \LaTeX file and code files to the Gradescope assignment `HW4 - Supplemental.' The supplemental files will only be checked in special cases, e.g. honor code issues, etc. Your files should be named in the same way as we provide them in the repository, e.g. \texttt{hw4.pdf}, etc.




\newpage

%%%%%%%%%%%%%%%%%%%%%%%%%%%%%%%%%%%%%%%%%%%%%
% Problem : Understanding the Transformer
%%%%%%%%%%%%%%%%%%%%%%%%%%%%%%%%%%%%%%%%%%%%%

\begin{problem}[Understanding the Transformer, 40 pts]

Transformers are the dominant architecture behind large language models and many modern
sequence processing systems.  At the heart of every transformer is \textbf{self-attention},
a mechanism that allows each position in a sequence to ``look at'' every other position
and decide how much information to gather from each.

In this problem you will (a)~work through the math of scaled dot-product attention,
(b)~understand key design choices, and (c)~implement the attention layer.

\medskip\noindent
\textbf{Background.}
Given an input matrix $\boldX \in \reals^{T \times d}$ (one row per token),
self-attention computes:
\[
  \text{Attention}(\boldQ, \boldK, \boldV)
  = \text{softmax}\!\left(\frac{\boldQ \boldK^\trans}{\sqrt{d_k}}\right)\boldV,
\]
where $\boldQ = \boldX \boldW_Q$, $\boldK = \boldX \boldW_K$,
$\boldV = \boldX \boldW_V$ with learnable weight matrices
$\boldW_Q, \boldW_K \in \reals^{d \times d_k}$ and
$\boldW_V \in \reals^{d \times d_v}$.
The softmax is applied row-wise: for a row vector $\boldz \in \reals^T$,
$\text{softmax}(\boldz)_i = e^{z_i} / \sum_{j=1}^{T} e^{z_j}$.

\begin{enumerate}

%%% Part (a) %%%
\item \textbf{[10 pts] Attention by hand.}

Consider a sequence of $T = 2$ tokens with embedding dimension $d = 2$. The input
matrix, query weight matrix, key weight matrix, and value weight matrix are:
\[
  \boldX = \begin{pmatrix} 1 & 0 \\ 0 & 1 \end{pmatrix}, \quad
  \boldW_Q = \begin{pmatrix} 1 & 1 \\ 0 & 1 \end{pmatrix}, \quad
  \boldW_K = \begin{pmatrix} 1 & 0 \\ 1 & 1 \end{pmatrix}, \quad
  \boldW_V = \begin{pmatrix} 2 & 0 \\ 0 & 1 \end{pmatrix}.
\]
So $d_k = d_v = 2$.

\begin{enumerate}
  \item Compute $\boldQ$, $\boldK$, $\boldV$, and the raw score matrix
  $\boldS = \boldQ \boldK^\trans$.

  \item Compute the scaled score matrix $\boldS / \sqrt{d_k}$.
  Then apply softmax row-wise to obtain the attention weight matrix
  $\boldA \in \reals^{2 \times 2}$.

  You may leave your answer in terms of $e^{(\cdot)}$ and sums, or compute
  decimal approximations.
\end{enumerate}
%%% Part (b) %%%
\item \textbf{[5 pts] Why scale by $\sqrt{d_k}$\,?}

Suppose the entries of $\boldq, \boldk \in \reals^{d_k}$ are independent random
variables with mean~$0$ and variance~$1$.

\begin{enumerate}
  \item Compute the mean and variance of the dot product
  $\boldq^\trans \boldk = \sum_{i=1}^{d_k} q_i k_i$.

  \item As $d_k$ grows, what happens to the magnitude of $\boldq^\trans \boldk$?
  Explain in one or two sentences what effect this has on the softmax distribution
  (i.e.\ does it become more uniform or more peaked?).

  \item Show that after scaling by $1/\sqrt{d_k}$ the dot product has variance~$1$
  regardless of $d_k$.
\end{enumerate}

%%% Part (c) %%%
\item \textbf{[5 pts] Permutation equivariance.}

Let $\boldP_\sigma \in \reals^{T \times T}$ be the permutation matrix corresponding
to a permutation $\sigma$ of $\{1, \dots, T\}$ (i.e.\ $(\boldP_\sigma)_{ij} =
\mathbf{1}[j = \sigma(i)]$).

\begin{enumerate}
  \item Show that self-attention (without positional encoding) is
  \textbf{permutation equivariant}: if we permute the input rows by replacing
  $\boldX$ with $\boldP_\sigma \boldX$, the output is permuted in the same way.
  That is, show
  \[
    \text{Attention}(\boldP_\sigma \boldX) = \boldP_\sigma \, \text{Attention}(\boldX).
  \]
  \textit{Hint}: What happens to $\boldQ$, $\boldK$, $\boldV$, and the softmax
  when you left-multiply $\boldX$ by $\boldP_\sigma$?

  \item Explain in two or three sentences why permutation equivariance is a
  problem for tasks like language modeling, where word order matters.

  \item The standard fix is to add a \textbf{positional encoding}
  $\boldP\boldE \in \reals^{T \times d}$ to the input:
  $\boldX' = \boldX + \boldP\boldE$.
  Why does this break permutation equivariance?
\end{enumerate}

%%% Part (d) %%%
\item \textbf{[10 pts] Implementing self-attention from scratch (code).}

In the accompanying notebook, implement single-head scaled dot-product self-attention
using \textbf{only NumPy} (do not use pre-existing implementations such as PyTorch and sklearn).

Your function should take:
\begin{itemize}
  \item $\boldX \in \reals^{T \times d}$: input embeddings
  \item $\boldW_Q \in \reals^{d \times d_k}$, $\boldW_K \in \reals^{d \times d_k}$,
        $\boldW_V \in \reals^{d \times d_v}$: weight matrices
\end{itemize}
and return:
\begin{itemize}
  \item $\text{output} \in \reals^{T \times d_v}$: the attention output
  \item $\text{weights} \in \reals^{T \times T}$: the attention weight matrix $\boldA$
\end{itemize}

\textbf{Implementation tips}:
\begin{itemize}
  \item For numerical stability, before applying softmax to a row~$\boldz$, subtract
  $\max(\boldz)$: compute $\text{softmax}(\boldz - \max(\boldz))$. (As an aside, note that this does not change the softmax function.)
  % \item Test your implementation against the provided test cases.
\end{itemize}

%%% Part (e) %%%
\item \textbf{[10 pts] Multi-head attention}

In the notebook, extend your implementation to \textbf{multi-head attention}
and use it inside a small transformer classifier trained on a synthetic
sequence dataset.

Given $h$ heads, each head
  uses its own $\boldW_Q^{(i)}, \boldW_K^{(i)} \in \reals^{d \times d_k}$ and
  $\boldW_V^{(i)} \in \reals^{d \times d_v}$, where $d_k = d_v = d / h$.
  The outputs of all heads are concatenated and projected:
  \[
    \text{MultiHead}(\boldX) = [\text{head}_1 \,;\, \cdots \,;\, \text{head}_h]\,
    \boldW_O,
  \]
  where $\boldW_O \in \reals^{d \times d}$.

After training, attach one learned attention matrix as a heatmap and
briefly describe what it suggests the model is attending to.


\end{enumerate}

\end{problem}


\newpage


\begin{solution}
	Your solution here.
\end{solution}


\newpage



\begin{problem}[Autoencoders, 76 pts]
Suppose we want to compress our data into a lower-dimensional representation. What is a principled way to do this with machine learning, such that we can be confident the compression retains as much useful information as possible?

\begin{enumerate}
        \item The Autoencoder. 
        
We can turn this into a game between two networks. An \textit{encoder} $f_\phi : \mathbb R^n \to \mathbb R^d$ takes an input $\mathbf x$ and compresses it down to $d$ dimensions: $\mathbf z = f_\phi(\mathbf x)$. A \textit{decoder} $g_\theta : \mathbb R^d \to \mathbb R^n$ must then reconstruct the original input from just this compressed representation: $\hat{\mathbf x} = g_\theta(\mathbf z)$. We measure how well the decoder succeeds by comparing $\hat{\mathbf x}$ to $\mathbf x$, and backpropagate through both networks:
\[
  \mathcal{L}_{\mathrm{recon}} = \frac{1}{N}\sum_{i=1}^{N} \|\mathbf x_i - g_\theta(f_\phi(\mathbf x_i))\|^2.
\]
This is the \textit{reconstruction loss}, and it is the core of how autoencoders are trained. The logic is simple: if the decoder can faithfully recover $\mathbf x$ from $\mathbf z$, then $\mathbf z$ must have captured the essential structure of $\mathbf x$.

Note that what makes a model an ``autoencoder'' is this encoder--decoder coupling trained with a reconstruction loss---not any particular architecture. The encoder and decoder can be MLPs, CNNs, Transformers, or anything else. It is the training paradigm that defines the autoencoder.

\begin{enumerate}

\item \textbf{Implementation}[12 pts]. In the notebook, implement a convolutional autoencoder on CelebA. Train it, and provide your training and test loss curves. Display the first 5 test images alongside their reconstructions.

\item \textbf{Sampling}[4 pts]. In the notebook, sample random vectors $\mathbf z \sim \mathcal{N}(\bm{0}, \mathbf{I})$ from the autoencoder's latent space and decode them. Do the outputs look like faces? In 1--2 sentences, explain why the autoencoder has no reason to make its latent space produce coherent outputs when sampled this way.
\end{enumerate}
\item The autoencoder gives us a compressed representation, but its latent space has no structure that would let us \textit{generate} new data. What if, instead of learning point embeddings, we tried to learn the probability distribution over images, and then sample from it?

The natural objective is maximum likelihood: given a generative model with latent variables $\bz$ and a prior $p(\bz) = \mathcal{N}(\bm{0}, \mathbf{I})$, the log-likelihood of a datapoint $\bx$ is
\[
  \log p(\bx) = \log \int p(\bx, \bz)\, d\bz = \log \int p(\bx \mid \bz)\, p(\bz)\, d\bz.
\]

To train this model via maximum likelihood, we would need the true posterior $p(\bz \mid \bx)$. By Bayes' theorem,
\[
  p(\bz \mid \bx) = \frac{p(\bx \mid \bz)\, p(\bz)}{p(\bx)}.
\]
However, computing $p(\bx)$ (the evidence) requires evaluating the integral above over all possible values of $\bz$, which is intractable when $p(\bx \mid \bz)$ is parameterized by a neural network in high dimensions.

\begin{enumerate}

\item \textbf{Intractability}[4 pts]. In 1--2 sentences, explain concretely why the integral $\int p(\bx \mid \bz)\,p(\bz)\,d\bz$ is intractable. Why can't we just estimate it by drawing samples $\bz^{(k)} \sim p(\bz)$ and averaging $p(\bx \mid \bz^{(k)})$?
\end{enumerate}

Since we cannot compute $p(\bz \mid \bx)$ exactly, we introduce a neural network $q_\phi(\bz \mid \bx)$ that serves as an \textit{approximate posterior}---our best learned guess at the true posterior. To measure how well $q_\phi(\bz \mid \bx)$ approximates $p(\bz \mid \bx)$, we use the Kullback--Leibler (KL) divergence and once again need a clever technique:
\[
  \KL\!\big(q_\phi(\bz \mid \bx) \,\|\, p(\bz \mid \bx)\big) = \E_{q_\phi(\bz \mid \bx)}\!\Big[\log q_\phi(\bz \mid \bx) - \log p(\bz \mid \bx)\Big].
\]

\begin{enumerate}
\setcounter{enumii}{1}
\item \textbf{Deriving the ELBO.} Starting from the KL divergence above:

\begin{enumerate}
  \item \textbf{Expanding} [8 pts]. Expand $\log p(\bz \mid \bx)$ using Bayes' rule, and rearrange to show:
  \[
    \log p(\bx) = \KL\!\big(q_\phi(\bz \mid \bx) \,\|\, p(\bz \mid \bx)\big) + \E_{q_\phi(\bz \mid \bx)}\!\Big[\log p(\bx, \bz) - \log q_\phi(\bz \mid \bx)\Big].
  \]

  \item \textbf{Interpretation} [4 pts]. We define the second term to be the \textbf{Evidence Lower Bound (ELBO)}:
  \[
    \mathrm{ELBO} = \E_{q_\phi(\bz \mid \bx)}\!\Big[\log p(\bx, \bz) - \log q_\phi(\bz \mid \bx)\Big].
  \]
  Since $\KL \geq 0$, conclude that $\mathrm{ELBO} \leq \log p(\bx)$. In one sentence, explain why this means maximizing the ELBO is a sensible surrogate for maximizing the intractable log-likelihood.

  \item\textbf{Rewriting} [8 pts]. Now show that the ELBO can be rewritten as:
  \[
    \mathrm{ELBO} = \E_{q_\phi(\bz \mid \bx)}\!\big[\log p(\bx \mid \bz)\big] - \KL\!\big(q_\phi(\bz \mid \bx) \,\|\, p(\bz)\big).
  \]
  (Hint: split $\log p(\bx, \bz) = \log p(\bx \mid \bz) + \log p(\bz)$.)

\end{enumerate}

\item \textbf{Interpreting the VAE loss.} The loss we minimize is the negative ELBO:
\[
  \mathcal{L}_{\mathrm{VAE}} = -\E_{q_\phi(\bz \mid \bx)}\!\big[\log p(\bx \mid \bz)\big] + \KL\!\big(q_\phi(\bz \mid \bx) \,\|\, p(\bz)\big).
\]

\begin{enumerate}
  \item \textbf{Interpretation}[4 pts]. The first term is the (negative) expected log-likelihood of $\bx$ given samples from the encoder. In 1--2 sentences, explain how this connects back to the reconstruction loss from Section 1. What does minimizing it encourage?

  \item \textbf{Interpretation}[4 pts]. The second term penalizes the encoder's output distribution $q_\phi(\bz \mid \bx)$ for deviating from the prior $p(\bz) = \mathcal{N}(\bm{0}, \mathbf{I})$. Explain how this term addresses the problem you identified in question 1(b): why does it give the model an incentive to organize the latent space for generation?

  \item \textbf{Tensity}[4 pts]. These two terms pull in opposite directions. In 1--2 sentences, describe the qualitative effect of this tension on the VAE's reconstructions and generated samples, compared to the standard autoencoder.
\end{enumerate}

\end{enumerate}
\item VAE Implementation
%% ============================================================

In a VAE, the encoder outputs the parameters of the approximate posterior\\ $q_\phi(\bz \mid \bx) = \mathcal{N}(\bm{\mu}_\phi(\bx),\, \mathrm{diag}(\bm{\sigma}_\phi^2(\bx)))$. That is, a shared convolutional backbone produces a feature vector, which is then projected through two separate linear heads: one outputting $\bm{\mu} \in \R^d$ and one outputting $\log \bm{\sigma}^2 \in \R^d$. During training, we sample $\bz$ from this distribution; during generation, we sample $\bz \sim \mathcal{N}(\bm{0}, \mathbf{I})$ and decode.

\begin{enumerate}

\item \textbf{Reparameterization}[8 pts]. We need to backpropagate through the sampling step $\bz \sim q_\phi(\bz \mid \bx)$, but sampling is not differentiable.
 In 1--2 sentences, explain why we cannot directly backpropagate gradients through a sampling operation with respect to the distribution's parameters.
  Show that if $\bm{\epsilon} \sim \mathcal{N}(\bm{0}, \mathbf{I})$, then $\bz = \bm{\mu}_\phi(\bx) + \bm{\sigma}_\phi(\bx) \odot \bm{\epsilon}$ has distribution $q_\phi(\bz \mid \bx)$, and explain briefly why this makes the loss differentiable w.r.t.\ $\phi$.

\item \textbf{Closed-form KL}[6 pts]. Derive the closed-form expression for the KL term when $q = \mathcal{N}(\bm{\mu}, \mathrm{diag}(\bm{\sigma}^2))$ and $p = \mathcal{N}(\bm{0}, \mathbf{I})$:
\[
  \KL(q \| p) = \frac{1}{2}\sum_{j=1}^{d}\Big(\mu_j^2 + \sigma_j^2 - \ln \sigma_j^2 - 1\Big).
\]
(Hint: work coordinate-by-coordinate. Use $\E[z_j] = \mu_j$ and $\E[z_j^2] = \mu_j^2 + \sigma_j^2$.)

\item \textbf{Implementation}[12 pts]. In the notebook, implement a VAE using the architecture described above. Train it, and provide your training and test loss curves (decomposed into reconstruction and KL terms). Display the first 5 test images alongside their reconstructions. Then, sample from $\mathcal{N}(\bm{0}, \mathbf{I})$ and decode---do the samples look like faces? Compare to your results from 1(b).

\end{enumerate}
\end{enumerate}
\end{problem}
\begin{solution}
	Your solution here.
\end{solution}
\newpage

\begin{problem}[Decision Trees, Random Forests and Mixture of Experts]

\noindent
\begin{enumerate}
	\item Suppose we have $B$ independently trained binary classifiers (e.g., decision trees), each with the same accuracy $p > \tfrac{1}{2}$ on a given test point. The ensemble predicts by \textbf{majority vote}.

Let $X$ be the number of trees that predict correctly. Then $X \sim \text{Binomial}(B, p)$, and the ensemble is correct whenever $X > B/2$.

\medskip
\noindent
\textbf{Background: Hoeffding's Inequality.}\quad
Let $Z_1, \ldots, Z_B$ be independent random variables with $Z_i \in [a_i, b_i]$. Hoeffding's inequality states that for the sample mean $\bar{Z} = \frac{1}{B}\sum_{i=1}^{B} Z_i$:
\[
P\!\big(\bar{Z} - \mathbb{E}[\bar{Z}] \leq -t\big) \;\leq\; \exp\!\!\left(\frac{-2B^2 t^2}{\sum_{i=1}^{B}(b_i - a_i)^2}\right).
\]
Intuitively, Hoeffding's inequality is a \emph{concentration inequality}: it says that the average of many bounded independent random variables is unlikely to be far from its expectation, and this probability shrinks \emph{exponentially} in the number of samples $B$. In our setting, each $Z_i \in \{0, 1\}$ indicates whether tree $i$ is correct, so $\bar{Z} = X/B$ and $\mathbb{E}[\bar{Z}] = p$.

\begin{enumerate}
    \item \textbf{[5 points]} Apply Hoeffding's inequality to show that the ensemble error rate satisfies
    \[
    P(X \leq \tfrac{B}{2}) \leq \exp\!\Big({-2B\big(p - \tfrac{1}{2}\big)^2}\Big).
    \]
    \emph{Hint:} The ensemble errs when $\bar{Z} \leq \tfrac{1}{2}$. 

    \item \textbf{[5 points]} Compute the bound for $p = 0.6$ and $B \in \{10, 100, 1000\}$. At what $B$ does the bound first drop below $10^{-6}$? Is that good or necessary?

    \item \textbf{[5 points]} The proof above assumes the trees are \emph{independent}. In practice, bagged trees are trained on overlapping bootstrap samples and are therefore correlated.
    \begin{enumerate}
        \item Explain intuitively why correlation weakens the convergence guarantee.
        \item Random forests add \textbf{feature subsampling} (each split considers only $m$ of $d$ features at random). Explain how this reduces correlation and partially restores convergence.
    \end{enumerate}
\end{enumerate}
\item
Part 1(c) showed that correlation between trees matters. We now make this precise. Suppose we have $B$ trees $T_1, \ldots, T_B$, each with the same prediction variance $\sigma^2 = \text{Var}(T_b(\mathbf{x}))$, and each pair has the same correlation $\rho = \text{Corr}(T_i(\mathbf{x}), T_j(\mathbf{x}))$ for $i \neq j$.

\begin{enumerate}
    \item \textbf{[8 points]} Prove that the variance of the ensemble prediction is:
    \[
    \text{Var}\!\left(\frac{1}{B}\sum_{b=1}^{B} T_b(\mathbf{x})\right) = \rho\,\sigma^2 + \frac{1-\rho}{B}\,\sigma^2.
    \]
    \emph{Hint:} Expand the variance of the sum, separate the diagonal terms ($i = j$) from the off-diagonal terms ($i \neq j$), and use $\text{Cov}(T_i, T_j) = \rho\,\sigma^2$.

    \item \textbf{[4 points]} Analyze the formula:
    \begin{enumerate}
        \item What does the ensemble variance converge to as $B \to \infty$? Can we make the variance arbitrarily small by adding more trees?
        \item What happens when $\rho = 0$? When $\rho = 1$? Relate each case to a practical scenario.
    \end{enumerate}

    \item \textbf{[4 points]} Random forests use $m = \lfloor\sqrt{d}\rfloor$ features per split (for classification). Using the formula from (d), explain why $m = 1$ (one random feature per split) is generally a bad idea, even though it minimizes $\rho$. What tradeoff does $m$ control?
\end{enumerate}

%% ============================================================
\item %% ============================================================
\medskip
\noindent
\textbf{Mixture of Experts (MoE).}\quad
A random forest is a \emph{dense ensemble}: every tree processes every input. A Mixture of Experts is a \emph{sparse ensemble}: a learned gating network routes each input to only a few specialized sub-networks (``experts''), and the output is a weighted combination:
\[
y = \sum_{k=1}^{K} g_k(\mathbf{x})\, E_k(\mathbf{x}), \quad \text{with most } g_k = 0 \text{ (sparse routing)}.
\]
This is the architecture behind many modern LLMs. Sparse MoE allows models to have many more total parameters (and thus more capacity) without proportionally increasing inference cost --- for example, Mixtral 8x7B has ${\sim}47$B total parameters but activates only ${\sim}13$B per token by routing to 2 of 8 experts. Unlike random forest trees, MoE experts \emph{specialize}: the gating network learns to send different types of inputs to different experts.

A key challenge in MoE training is \textbf{expert collapse}, where the gating network funnels most inputs to one or two experts while the rest go unused, wasting capacity. This is typically addressed with an auxiliary \emph{load-balancing loss} that penalizes uneven routing.

\textbf{[4 points]} In 3--5 sentences, compare how random forests and MoE each handle the tradeoff between model capacity and computational cost. Why can a random forest not ``scale up'' in the way that MoE can?
\end{enumerate}

\end{problem}

\begin{solution}
	Your solution here.
\end{solution}


\newpage 




\textbf{Name}:

\textbf{Collaborators and Resources}: 

\end{document}
